\let\negmedspace\undefined
\let\negthickspace\undefined
\documentclass[journal]{IEEEtran}
\usepackage[a5paper, margin=10mm, onecolumn]{geometry}
%\usepackage{lmodern} % Ensure lmodern is loaded for pdflatex
\usepackage{tfrupee} % Include tfrupee package

\setlength{\headheight}{1cm} % Set the height of the header box
\setlength{\headsep}{0mm}     % Set the distance between the header box and the top of the text

\usepackage{cite}
\usepackage{graphicx}
\usepackage{amsmath,amssymb,amsfonts,amsthm}
\usepackage{algorithmic}
\usepackage{graphicx}
\usepackage{textcomp}
\usepackage{xcolor}
\usepackage{txfonts}
\usepackage{listings}
\usepackage{enumitem}
\usepackage{mathtools}
\usepackage{gensymb}
\usepackage{comment}
\usepackage[breaklinks=true]{hyperref}
\usepackage{tkz-euclide} 
\usepackage{listings}
\usepackage{gvv}                                        
%\def\inputGnumericTable{}                                 
\usepackage[latin1]{inputenc} 
\usetikzlibrary{arrows.meta, positioning}
\usepackage{xparse}
\usepackage{color}                                            
\usepackage{array}                                            
\usepackage{longtable}                                       
\usepackage{calc}                                             
\usepackage{multirow}
\usepackage{multicol}
\usepackage{hhline}                                           
\usepackage{ifthen}                                           
\usepackage{lscape}
\usepackage{tabularx}
\usepackage{array}
\usepackage{float}
\newtheorem{theorem}{Theorem}[section]
\newtheorem{problem}{Problem}
\newtheorem{proposition}{Proposition}[section]
\newtheorem{lemma}{Lemma}[section]
\newtheorem{corollary}[theorem]{Corollary}
\newtheorem{example}{Example}[section]
\newtheorem{definition}[problem]{Definition}
\newcommand{\BEQA}{\begin{eqnarray}}
\newcommand{\EEQA}{\end{eqnarray}}
\usepackage{float}
\theoremstyle{remark}
\usepackage{circuitikz}
\usepackage{tikz} 
\graphicspath{{figs/}}

\begin{document}
\title{4.3.6}
\author{EE25BTECH11022 - sankeerthan}
% \maketitle
% \newpage
% \bigskip
{\let\newpage\relax\maketitle}

\renewcommand{\thefigure}{\theenumi}
\renewcommand{\thetable}{\theenumi}
\setlength{\intextsep}{10pt} % Space between text and floats

\textbf{Question}:passing through the points \brak{3,4,-7} and \brak{1,-1,6}.

\textbf{solution}:let Passing through the points $\vec{A} =\myvec{3 \\4 \\ -7}$ and $\vec{B} =\myvec{1 \\ -1 \\ 6}$. \\
\myvec{3 & 4 & -7 \\ 1 & -1 & 6}$\mathbf{n}$ =\myvec{1 \\ 1} \\
Writing the augmented matrix:
\begin{align}
\myvec{3 & 4 & -7 & 1 \\ 1 & -1 & 6 & 1} \xleftrightarrow{R_1 \leftarrow R_1 - 3R_2} \myvec{0 & 7 & -25 & -2 \\ 1 & -1 & 6 & 1}
\xleftrightarrow{R_1 \leftrightarrow R_2}\myvec{1 & -1 & 6 & 1 \\ 0 & 7 & -25 & -2}
\end{align}
\begin{align}
n_1 - n_2 + 6 n_3 = 1 \\
7 n_2 - 25 n_3 = -2 
\end{align}
From third equation: 
\begin{align}
n_2 = \frac{-2 + 25 n_3}{7}
\end{align}
Substitute in second equation: 
\begin{align}
n_1 = \frac{5 - 17 n_3}{7} 
\end{align}
Let $n_3 = t$ , then
\begin{align}
\mathbf{n} = \myvec{\frac{5 - 17 t}{7} \\ \frac{-2 + 25 t}{7} \\ t}
\end{align}
let $t = 1$, multiply with $7$ to clear denominators:
\begin{align}
\mathbf{n} = \myvec{-12 \\ 23 \\ 7}
\end{align}

Thus, the equation of the line is:
$\myvec{-12 & 23 & 7}\mathbf{x}$ = 7

\begin{figure}
   \centering
   \includegraphics[width=1\columnwidth]{figs/points.png}
   \caption{}
   \label{}
\end{figure}
\end{document}